../cictikz/src/cictikz/data/tex/cictikz_preamble.tex

% The D flip-flop is two latches back to back: the master is
% transparent while the slave holds, and the clock inverter swaps them.

\begin{circuitikz}[thick]
  ../cictikz/src/cictikz/data/tex/cictikz_lib.tex

  % ----------------------------------------------------------------
  % Flip-flop symbol
  % ----------------------------------------------------------------
  \draw (-8.2,-1.4) rectangle (-5.8,1.6);
  \node[anchor=west] at (-8.15,1.1)  {$D$};
  \node[anchor=east] at (-5.85,1.1)  {$Q$};
  \node[anchor=east] at (-5.85,-0.6) {$\overline{Q}$};
  \draw (-7.6,-1.4) -- (-7.3,-1.15) -- (-7.0,-1.4);
  \draw (-7.3,-1.4) -- (-7.3,-2.0);
  \draw (-8.8,1.1) -- (-8.2,1.1);
  \draw (-5.8,1.1) -- (-5.2,1.1);
  \draw (-5.8,-0.6) -- (-5.2,-0.6);

  \draw[armygreen, very thick, ->] (-4.4,0.1) -- (-3.2,0.1);

  % ----------------------------------------------------------------
  % Master and slave latches
  % ----------------------------------------------------------------

  % master
  \draw (0,-1.4) rectangle (2.4,1.6);
  \node[anchor=west] at (0.05,1.1)  {$D$};
  \node[anchor=east] at (2.35,1.1)  {$Q$};
  \node[anchor=east] at (2.35,-0.6) {$\overline{Q}$};
  \node[scale=1.4] at (1.1,0.3) {$L$};
  \draw (0.6,-1.4) -- (0.9,-1.15) -- (1.2,-1.4);
  \draw (2.4,-0.6) -- (3.0,-0.6);

  % slave
  \draw (4.6,-1.4) rectangle (7.0,1.6);
  \node[anchor=west] at (4.65,1.1)  {$D$};
  \node[anchor=east] at (6.95,1.1)  {$Q$};
  \node[anchor=east] at (6.95,-0.6) {$\overline{Q}$};
  \node[scale=1.4] at (5.7,0.3) {$L$};
  \draw (5.2,-1.4) -- (5.5,-1.15) -- (5.8,-1.4);

  % data path
  \draw (-1.3,1.1) node[anchor=east] {$D$} to [short,o-] (0,1.1);
  \draw (2.4,1.1) -- (4.6,1.1);
  \draw (7.0,1.1) to [short,-o] ++(0.7,0) node[anchor=west] {$Q$};
  \draw (7.0,-0.6) to [short,-o] ++(0.7,0) node[anchor=west] {$\overline{Q}$};

  % clock: straight into the master, inverted into the slave
  \draw (-1.3,-3.4) node[anchor=east] {$C$} to [short,o-] (0.9,-3.4);
  \fill (0.9,-3.4) circle (0.075);
  \draw (0.9,-3.4) -- (0.9,-1.4);
  \draw (2.35,-3.4) \cicInv{cinv};
  \draw (0.9,-3.4) -- (cinv_in);
  \draw (cinv_out) -- (5.5,-3.4) -- (5.5,-1.4);

\end{circuitikz}
\end{document}
